\documentclass[11pt]{article}
\usepackage[letterpaper,margin=0.7in]{geometry}
\usepackage{amsmath,amssymb}
\usepackage{booktabs}
\usepackage{array}
\usepackage{xcolor}
\usepackage{tcolorbox}
\usepackage{titlesec}
\usepackage{parskip}

\definecolor{garnet}{HTML}{73000A}
\definecolor{atlantic}{HTML}{466A9F}
\definecolor{congaree}{HTML}{1F414D}
\definecolor{horseshoe}{HTML}{65780B}
\definecolor{tenblack}{HTML}{ECECEC}

\tcbset{
  refbox/.style={
    colback=tenblack, colframe=garnet, sharp corners,
    boxrule=0.6pt, left=5pt, right=5pt, top=4pt, bottom=4pt
  }
}

\titleformat{\section}{\Large\bfseries\color{garnet}}{\thesection}{0.5em}{}[{\titlerule[0.6pt]}]
\titleformat{\subsection}{\large\bfseries\color{congaree}}{\thesubsection}{0.5em}{}

\renewcommand{\arraystretch}{1.25}

\title{\color{garnet}\textbf{Set \& Logic Notation Reference}}
\author{CSCE 763 --- J.C. Vaught}
\date{}

\begin{document}
\maketitle
\vspace{-2em}
\hrule
\vspace{1em}

\section{Set Operations (use these on the exam)}

\begin{tcolorbox}[refbox]
\begin{tabular}{@{}clll@{}}
\toprule
\textbf{Symbol} & \textbf{Name} & \textbf{Meaning} & \textbf{Plain English} \\
\midrule
$\cap$ & Intersection & $A \cap B$: elements in both & \textbf{AND} \\
$\cup$ & Union & $A \cup B$: elements in either & \textbf{OR} \\
$A - B$ & Difference & in $A$ but not $B$ & A AND NOT B \\
$A^c$ & Complement & not in $A$ & \textbf{NOT} \\
$\subseteq$ & Subset & every element of left is in right & ``fits inside'' \\
$\subset$ & Strict subset & subset and not equal & ``fits, smaller'' \\
$\in$ & Element of & $x \in A$ & ``x is in A'' \\
$\notin$ & Not element of & $x \notin A$ & ``x is not in A'' \\
$\emptyset$ & Empty set & no elements & ``nothing'' \\
$|A|$ & Cardinality & number of elements & ``size of A'' \\
\bottomrule
\end{tabular}
\end{tcolorbox}

\section{Morphology Operators}

\begin{tcolorbox}[refbox]
\begin{tabular}{@{}cll@{}}
\toprule
\textbf{Symbol} & \textbf{Name} & \textbf{Definition} \\
\midrule
$\hat{B}$ & Reflection & $\{w \mid w = -b,\, b \in B\}$ \\
$(B)_z$ & Translation & $\{c \mid c = b + z,\, b \in B\}$ \\
$\ominus$ & Erosion & $A \ominus B = \{z \mid (B)_z \subseteq A\}$ \\
$\oplus$ & Dilation & $A \oplus B = \{z \mid (\hat{B})_z \cap A \ne \emptyset\}$ \\
$\circ$ & Opening & $A \circ B = (A \ominus B) \oplus B$ \\
$\bullet$ & Closing & $A \bullet B = (A \oplus B) \ominus B$ \\
$\circledast$ & Hit-or-miss & $A \circledast B = (A \ominus B_1) \cap (A^c \ominus B_2)$ \\
\bottomrule
\end{tabular}
\end{tcolorbox}

\section{Pure Logic (propositions, not sets)}

\begin{tcolorbox}[refbox]
\begin{tabular}{@{}cll@{}}
\toprule
\textbf{Symbol} & \textbf{Name} & \textbf{Meaning} \\
\midrule
$\wedge$ & AND & both true \\
$\vee$ & OR & at least one true \\
$\neg$ or $\sim$ & NOT & flip truth \\
$\Rightarrow$ & Implies & if-then \\
$\Leftrightarrow$ & If and only if & both directions \\
$\forall$ & For all & every \\
$\exists$ & There exists & at least one \\
\bottomrule
\end{tabular}
\end{tcolorbox}

\textbf{Key distinction.} $\wedge$/$\vee$ work on \emph{propositions} (true/false statements). $\cap$/$\cup$ work on \emph{sets}. They mean the same thing logically but appear in different contexts. On a morphology or segmentation exam, you almost always use $\cap$ and $\cup$.

\section{How These Show Up in Course Topics}

\subsection{Region Segmentation}
\begin{itemize}
\item Disjoint regions: $R_i \cap R_j = \emptyset$ for $i \ne j$ (no overlap).
\item Cover the image: $\bigcup_{i=1}^{n} R_i = R$ (every pixel is in some region).
\item Merge predicate: if $Q(R_i \cup R_j) = \text{TRUE}$, merge $R_i$ and $R_j$.
\item Each region must be connected: every pair of pixels in $R_i$ has a path entirely inside $R_i$.
\end{itemize}

\subsection{Connectivity / Components}
\begin{itemize}
\item Pixels $p, q$ are connected in $S$ if a path of pixels in $S$ joins them.
\item Connected component $C \subseteq S$: maximal connected subset.
\item $S = C_1 \cup C_2 \cup \cdots \cup C_k$ with $C_i \cap C_j = \emptyset$ for $i \ne j$.
\end{itemize}

\subsection{Morphology Examples}
\begin{itemize}
\item Erosion ``fits'': SE must fit entirely inside $A$.
\item Dilation ``hits'': reflected SE must overlap $A$ at least once.
\item Boundary: $\beta(A) = A - (A \ominus B)$.
\item Duality: $A \oplus B = (A^c \ominus \hat{B})^c$.
\item Idempotence: $(A \circ B) \circ B = A \circ B$, $(A \bullet B) \bullet B = A \bullet B$.
\end{itemize}

\section{Translation Quick-Test}

\begin{tcolorbox}[refbox]
\begin{tabular}{@{}ll@{}}
\toprule
\textbf{Expression} & \textbf{Plain English} \\
\midrule
$A \cap B = \emptyset$ & A and B do not overlap \\
$A \subseteq B$ & A fits inside B \\
$p \in (A \cap B)$ & pixel $p$ is in both regions \\
$A \cup B$ & combine A and B into one set \\
$A^c$ & the background (everything not in A) \\
$A - B$ & A with B removed \\
$|A| > 0$ & A is nonempty \\
$\forall p \in A : f(p) > T$ & every pixel in A has intensity above T \\
$\exists p \in R : Q(p) = \text{TRUE}$ & some pixel in R satisfies Q \\
\bottomrule
\end{tabular}
\end{tcolorbox}

\section{Five-Symbol Legend (worth putting on cheat sheet)}

\begin{center}
\fbox{\parbox{0.85\textwidth}{\centering
$\cap$ = AND / intersection \quad
$\cup$ = OR / union \quad
$\subseteq$ = subset \quad
$A^c$ = complement \quad
$\emptyset$ = empty set
}}
\end{center}

That covers essentially every set-theory symbol the final exam can throw at you.

\end{document}
